%=============================================================================
% WAVE 3, ITERATION 3: PATENT 1 UPDATE
% Field Drag / Cloaking --- Gravitational-Doppler Sensing (Precision),
% Kerr Analog Frame-Dragging, Atom Interferometer Detectability,
% Non-Reciprocal Drag, Cross-Patent P1/P5/P6 SNR
%
% Agent 1 | DFD Research Programme
% Based on: W3i2 P1 update, W3i2 P5/P6 updates,
%           DFD_Psi_Field_Drag_Reduction_Deep_Analysis.tex,
%           DFD Formalism, Master Corpus
% Date: 19 March 2026
%=============================================================================

\documentclass[11pt,twocolumn]{article}
\usepackage[utf8]{inputenc}
\usepackage{amsmath,amssymb,amsfonts,amsthm}
\usepackage{geometry}
\usepackage{booktabs}
\usepackage{siunitx}
\usepackage{hyperref}
\usepackage{xcolor}
\usepackage{enumitem}
\usepackage{mathtools}

\geometry{margin=1in}

\newcommand{\dpsi}{\delta\psi}
\newcommand{\lam}{\lambda}
\newcommand{\Opsi}{\Omega_\psi}
\newcommand{\gpsi}{\gamma_\psi}
\newcommand{\Mpsi}{M_\psi}
\newcommand{\astar}{a_\star}
\newcommand{\Rey}{\mathrm{Re}}
\newcommand{\vvec}{\mathbf{v}}
\newcommand{\nvec}{\hat{\mathbf{n}}}
\newcommand{\order}[1]{\mathcal{O}(#1)}
\newcommand{\half}{\tfrac{1}{2}}
\newcommand{\Ri}{\mathrm{Ri}}
\newcommand{\SNR}{\mathrm{SNR}}

\definecolor{errorcolor}{rgb}{0.7,0,0}
\definecolor{newcolor}{rgb}{0,0.4,0}
\definecolor{conjecturecolor}{rgb}{0.6,0.3,0}

\newtheorem{theorem}{Theorem}
\newtheorem{proposition}{Proposition}
\newtheorem{corollary}{Corollary}
\newtheorem{lemma}{Lemma}

\title{Wave 3, Iteration 3: Patent 1 Update\\
Field Drag and Cloaking ---\\
Gravitational-Doppler Sensing (Precision Analysis),\\
Kerr Analog Frame-Dragging, Atom Interferometer\\
Detectability, Non-Reciprocal Drag,\\
and Combined P1/P5/P6 Cross-Patent SNR}

\author{DFD Research Programme --- Agent 1\\
\textit{Patent 1: Field Drag, Hydrodynamic Cloaking, Refractive Bubble}}

\date{19 March 2026}

\begin{document}
\maketitle

%=============================================================================
\section*{W3i3: P1 Field Drag and Cloaking --- Iteration 3 Update}
%=============================================================================

\noindent This document is the third iteration of the Wave~3 cross-reference
cycle for Patent~1. W3i1 established the lambda-drag threshold and K-H
lever arm. W3i2 established: (1)~the $\mu_0$-corrected drag threshold
$|\lam-1|_{\rm drag} = 1.85\times10^{-2}$; (2)~K-H suppression five orders
below drag threshold via the $c^2$ lever arm; (3)~simultaneous drag reduction
and acoustic cloaking at $|\dpsi|\approx0.25$--$0.31$; (4)~the instability
sandwich for wall-concentrated psi-fields; (5)~the pivot from drag reduction
to gravitational-Doppler sensing as the near-term experimental direction.

Iteration~3 goes deeper on six topics not addressed in W3i2:
\begin{enumerate}[nosep]
\item Gravitational-Doppler formula precision: what frequency shift for a
  LEO satellite passing through a density gradient?
\item Kerr-analog effects: does DFD predict frame-dragging-like signatures
  from rotating masses distinct from GR?
\item Minimum detectable psi-gradient using existing atom interferometers
  (AION, MAGIS).
\item Non-reciprocal drag: is the force different for $+v$ vs $-v$ through
  a psi-gradient?
\item P1 + P6 quantum sensors: combined SNR calculation.
\item P1 + P5 timing: can timing anomalies near dense objects reveal
  psi-gradients?
\end{enumerate}

All derivations are carried out in full with explicit numbers.
Uncertain claims are marked
\textcolor{conjecturecolor}{[CONJECTURE]} or
\textcolor{errorcolor}{[NEEDS VERIFICATION]}.

%=============================================================================
\section*{W3i3 \S1: Gravitational-Doppler Sensing: Precision Analysis}
\label{sec:grav_doppler}
%=============================================================================

\subsection*{1.1 Clock Rate Formula in the DFD Metric}

The DFD optical metric at weak field, low velocity is:
\begin{equation}
ds^2 = -c^2 e^{-\dpsi}dt^2 + e^{+\dpsi}(dx^2+dy^2+dz^2).
\label{eq:metric}
\end{equation}

The proper time of an on-board atomic clock moving with coordinate velocity
$v \ll c$:
\begin{align}
\frac{d\tau}{dt} &= \sqrt{e^{-\dpsi} - e^{+\dpsi}\frac{v^2}{c^2}} \notag\\
&\approx e^{-\dpsi/2}\left(1 - \half e^{2\dpsi}\frac{v^2}{c^2}\right)^{1/2} \notag\\
&\approx 1 - \frac{\dpsi}{2} - \frac{v^2}{2c^2}
  + \order{\dpsi^2, \dpsi\, v^2/c^2}.
\label{eq:clock_rate}
\end{align}

The fractional frequency shift relative to a reference clock at spatial
infinity ($\dpsi_\infty = 0$, $v_{\rm ref} = 0$):
\begin{equation}
\frac{\Delta\nu}{\nu_0} = \frac{d\tau/dt - 1}{1}
\approx -\frac{\dpsi}{2} - \frac{v^2}{2c^2}.
\label{eq:freq_shift}
\end{equation}

The \textbf{psi-only contribution}:
\begin{equation}
\left.\frac{\Delta\nu}{\nu_0}\right|_\psi = -\frac{\dpsi}{2}.
\label{eq:psi_clock}
\end{equation}

Note: this differs by a factor of 2 from the photon gravitational redshift
($\Delta\nu/\nu = -\dpsi$ for a photon between two potential levels),
because the clock formula involves a \emph{single} evaluation of the potential,
whereas the photon redshift involves the difference between emission and
reception potential. This factor-of-2 distinction is testable and is
the same as in GR (geodetic vs. gravitational redshift).

\subsection*{1.2 Rate-of-Change Signal for Moving Satellite}

As the satellite traverses a psi-gradient, the clock rate drifts:
\begin{equation}
\frac{d}{dt}\!\left(\frac{\Delta\nu}{\nu_0}\right)_\psi
= -\half\frac{d\dpsi}{dt}
= -\half\,\vvec_{\rm sat}\cdot\nabla\dpsi.
\label{eq:freq_shift_rate}
\end{equation}

This is the primary observable: not the static frequency offset (buried in
GR), but the \emph{rate of change} correlated with satellite trajectory
through a known psi-source.

\subsection*{1.3 Circular LEO Orbit in Earth's Background Psi-Field}

In DFD the Newtonian limit gives $\dpsi = -2GM_\oplus/(c^2 r)$. The
gradient:
\begin{equation}
\frac{\partial\dpsi}{\partial r}\bigg|_{\rm LEO}
= \frac{2GM_\oplus}{c^2 r^2}
= \frac{2\times3.984\times10^{14}}{9\times10^{16}\times(6.771\times10^6)^2}
= 1.930\times10^{-16}~\text{m}^{-1}.
\label{eq:earth_gradient}
\end{equation}

For a \textit{circular} orbit the radial velocity is identically zero,
so $\vvec_{\rm sat}\cdot\nabla\dpsi = 0$: a circular orbit produces
\emph{no time-varying} psi frequency shift. The background signal vanishes.

\subsection*{1.4 Eccentric Orbit: Analytic Signal Formula}

For an elliptical orbit with semi-major axis $a$, eccentricity $e$, the
radial velocity is:
\begin{equation}
v_r = \frac{na e\sin E}{\sqrt{1-e^2}},
\qquad n = \sqrt{GM_\oplus/a^3}.
\label{eq:radial_vel}
\end{equation}

The psi-frequency-shift rate:
\begin{equation}
\frac{d}{dt}\!\left(\frac{\Delta\nu}{\nu_0}\right)_\psi
= -\half v_r\,\frac{\partial\dpsi}{\partial r}
= -\frac{G M_\oplus\, v_r}{c^2 r^2}.
\label{eq:ecc_signal}
\end{equation}

The peak-to-peak frequency excursion over one orbit (integrating
$\Delta\nu/\nu_0$ from periapsis to apoapsis):
\begin{equation}
\left|\Delta\!\left(\frac{\Delta\nu}{\nu_0}\right)_\psi\right|
= |\dpsi_{\rm apo} - \dpsi_{\rm peri}|/2
= \frac{GM_\oplus}{c^2}\left(\frac{1}{r_{\rm peri}} - \frac{1}{r_{\rm apo}}\right)
= \frac{GM_\oplus}{c^2}\cdot\frac{2e}{a(1-e^2)}.
\label{eq:orbit_excursion}
\end{equation}

\textbf{LEO, $e=0.1$}: $a = 6771$~km, $r_{\rm peri} = 6094$~km,
$r_{\rm apo} = 7448$~km.
\begin{equation}
\left|\Delta\!\left(\frac{\Delta\nu}{\nu_0}\right)_\psi\right|
= \frac{3.984\times10^{14}}{9\times10^{16}}
  \cdot\frac{0.2}{6.771\times10^6\times0.99}
= 4.43\times10^{-3}
  \times\frac{0.2}{6.703\times10^6}
= 1.32\times10^{-10}.
\label{eq:orbit_excursion_numerical}
\end{equation}

This is the total psi-channel excursion. However, this is again the
\textit{standard GR} gravitational redshift variation over the orbit ---
identical to what GR predicts, since DFD reproduces Newtonian potential
at leading order.

\subsection*{1.5 DFD-Specific (Non-GR) Frequency Signal}

The DFD deviation from GR arises at post-Newtonian order. The
leading correction to the clock rate:
\begin{align}
\frac{d\tau}{dt}\bigg|_{\rm DFD} - \frac{d\tau}{dt}\bigg|_{\rm GR}
&\approx -\frac{\dpsi^2}{8} + \frac{v^2\dpsi}{2c^2} + \order{v^4/c^4}
\label{eq:DFD_vs_GR}
\end{align}
\textcolor{conjecturecolor}{[CONJECTURE: this expansion follows from
expanding Eq.~\eqref{eq:clock_rate} to second order and comparing with
the standard PPN clock formula; requires verification against the full
DFD PPN parameter derivation.]}

At LEO: $|\dpsi_{\rm LEO}| = 1.306\times10^{-9}$, $v^2/c^2 = 6.53\times10^{-10}$:
\begin{equation}
\left|\frac{d\tau}{dt}\bigg|_{\rm DFD} - \frac{d\tau}{dt}\bigg|_{\rm GR}\right|
\approx \frac{(1.306\times10^{-9})^2}{8}
= 2.1\times10^{-19}.
\label{eq:DFD_GR_diff}
\end{equation}

Current GPS clock stability: $\sigma_y \approx 10^{-13}$ at 1~s, or
$\sim10^{-14}$ per day. The DFD-vs-GR deviation of $2\times10^{-19}$ is
five orders of magnitude below current GPS clock noise. \textbf{Conclusion:
DFD and GR are indistinguishable via LEO satellite clocks from the
background Earth field.}

\subsection*{1.6 Controlled Psi-Source: The Useful Experimental Configuration}

The physically meaningful signal is from a \emph{controlled EM source}
generating a psi-perturbation $\dpsi_{\rm EM}$ that the satellite then
observes. For a Gaussian feature of amplitude $A$ and width $\ell$:
\begin{equation}
\dpsi_{\rm EM}(x) = A\exp\!\left(-\frac{x^2}{2\ell^2}\right),
\label{eq:gaussian_feature}
\end{equation}
where $x = v_{\rm sat}(t - t_0)$ is the along-track displacement. The
satellite clock frequency variation:
\begin{equation}
\frac{d}{dt}\!\left(\frac{\Delta\nu}{\nu_0}\right)_{\rm EM}
= \frac{v_{\rm sat}Ax}{2\ell^2}
  \exp\!\left(-\frac{x^2}{2\ell^2}\right).
\label{eq:gaussian_signal}
\end{equation}

Peak frequency-shift rate: at $x = \ell$:
\begin{equation}
\left.\frac{d}{dt}\!\left(\frac{\Delta\nu}{\nu_0}\right)_{\rm EM}\right|_{\rm peak}
= \frac{v_{\rm sat}|A|}{2\ell\sqrt{e}}.
\label{eq:peak_rate}
\end{equation}

The peak-to-peak frequency excursion (from $x=-\infty$ to $+\infty$):
\begin{equation}
\Delta\!\left(\frac{\Delta\nu}{\nu_0}\right)_{\rm EM} = |A|/2.
\label{eq:pp_excursion}
\end{equation}

\textbf{Detectability criterion}: GPS clock stability $\sigma_y(1\,\text{s})
\approx 10^{-13}$. Signal detectable if $|A|/2 > 10^{-13}$, i.e.,
$|A| > 2\times10^{-13}$.

From W3i2 Eq.~(10) with $u_{\rm EM} = P/(c\pi w_0^2)$ and P = 100~kW,
$w_0 = 0.05$~m at the source:
\begin{equation}
A \approx (\lam-1)\times2.8\times10^{-22}.
\label{eq:A_from_lambda}
\end{equation}

For detection: $|\lam-1| > 2\times10^{-13}/(2.8\times10^{-22})
= 7.1\times10^{8}$.

\textbf{Negative finding}: Detecting a 100~kW EM source psi-signal via
LEO satellite clock anomaly requires $|\lam-1| > 7\times10^8$ --- nine
orders of magnitude above the current passive bound.

\subsubsection*{Gap summary for Section 1}

\begin{center}
\begin{tabular}{@{}lcc@{}}
\toprule
\textbf{Configuration} & Required $|\lam-1|$ & Gap to bound \\
\midrule
Earth background (GR deviation) & Undetectable ($\sim10^{12}$) & $10^{12}$ \\
100~kW source, GPS clock & $7\times10^8$ & $1.4\times10^{15}$ \\
100~kW source, optical clock & $7\times10^5$ & $1.4\times10^{12}$ \\
\bottomrule
\end{tabular}
\end{center}

\textit{Note}: Optical clocks achieve $\sigma_y \sim 10^{-18}$, improving
the required $|\lam-1|$ to $7\times10^5$ --- still $10^{12}$ above the
current bound.

%=============================================================================
\section*{W3i3 \S2: Kerr Analog Effects}
\label{sec:kerr}
%=============================================================================

\subsection*{2.1 Frame-Dragging in GR and DFD}

In linearized GR, a rotating mass generates a gravitomagnetic field:
\begin{equation}
\mathbf{B}_{\rm GR} = \frac{G}{c^2}\left[
  \frac{3(\mathbf{J}\cdot\hat r)\hat r - \mathbf{J}}{r^3}
\right],
\label{eq:B_GR}
\end{equation}
which causes gyroscope precession (Lense-Thirring effect) at rate
$\dot\Omega_{\rm LT} = \mathbf{B}_{\rm GR}/2$.

In DFD, the gravitomagnetic field arises from the same mass-current source
through the $g_{0i}$ terms in the DFD metric, which at leading post-Newtonian
order must reproduce GR (by construction, since DFD is required to pass all
Solar System tests). The DFD \emph{deviation} from GR frame-dragging arises
only from the additional EM source term $(\lam-1)\mathcal{L}_{\rm EM}$.

\subsection*{2.2 EM-Sourced DFD Gravitomagnetic Correction}

The additional DFD psi-field from a rotating electromagnetically active
body. For a magnetic dipole $\mu_m$ rotating at angular velocity $\Omega$,
the EM energy density has a time-dependent component at frequency $2\Omega$
(due to the rotating quadrupole pattern). The time-averaged contribution
to the psi-field gradient at distance $r$:
\textcolor{conjecturecolor}{[CONJECTURE: the following derivation is
approximate; the full calculation requires the DFD tensor gravitomagnetic
equations which have not been derived in the available corpus.]}

\begin{equation}
\langle\nabla\dpsi_{\rm EM}\rangle
\approx \frac{(\lam-1)G\mu_0\mu_m^2\Omega}{4\pi c^6 r^5}
\times f(\theta),
\label{eq:gm_EM_correction}
\end{equation}
where $f(\theta)$ is an angular factor $\order{1}$.

The effective frame-dragging correction to gyroscope precession:
\begin{equation}
\delta\dot\phi_{\rm DFD} \approx c^2\langle\nabla\dpsi_{\rm EM}\rangle/2
= \frac{(\lam-1)G\mu_0\mu_m^2\Omega}{8\pi c^4 r^5}.
\label{eq:gm_precession}
\end{equation}

\subsection*{2.3 Numerical Evaluation for Earth}

Earth parameters: $\mu_m^\oplus = 8.0\times10^{22}$~A$\cdot$m$^2$,
$\Omega_\oplus = 7.27\times10^{-5}$~rad/s, $r = R_\oplus = 6.371\times10^6$~m,
$\mu_0 = 4\pi\times10^{-7}$~H/m.

\begin{align}
\delta\dot\phi_{\rm DFD}^\oplus
&= \frac{|\lam-1|\times 6.674\times10^{-11}
  \times 4\pi\times10^{-7}
  \times (8\times10^{22})^2
  \times 7.27\times10^{-5}}
  {8\pi\times (3\times10^8)^4\times (6.371\times10^6)^5}
\label{eq:earth_gm_calc}
\end{align}

Numerator: $6.674\times10^{-11}\times4\pi\times10^{-7}
\times6.4\times10^{45}\times7.27\times10^{-5}$
$= 6.674\times4\pi\times6.4\times7.27\times10^{-11-7+45-5}$
$= 6.674\times12.566\times46.53\times10^{22}$
$= 3904\times10^{22} = 3.904\times10^{25}$.

Denominator: $8\pi\times(3\times10^8)^4\times(6.371\times10^6)^5$
$= 25.133\times8.1\times10^{33}\times1.085\times10^{34}$
$= 25.133\times8.786\times10^{67}$
$= 2.208\times10^{69}$.

\begin{equation}
\delta\dot\phi_{\rm DFD}^\oplus
= |\lam-1|\times\frac{3.904\times10^{25}}{2.208\times10^{69}}
= |\lam-1|\times1.77\times10^{-44}~\text{rad/s}.
\label{eq:earth_gm_result}
\end{equation}

At current bound $|\lam-1| = 4.9\times10^{-7}$:
\begin{equation}
\delta\dot\phi_{\rm DFD}^\oplus = 4.9\times10^{-7}
\times1.77\times10^{-44}
= 8.7\times10^{-52}~\text{rad/s}.
\end{equation}

The GR Lense-Thirring precession at Earth's surface is
$\approx 6.3\times10^{-14}$~rad/s (confirmed by Gravity Probe~B to 1\%).

\begin{proposition}[DFD Kerr-analog is $10^{38}$ below GR frame-dragging]
The DFD EM-sourced gravitomagnetic correction to frame-dragging around
Earth is $8.7\times10^{-52}$~rad/s at the current passive bound,
which is $\sim 7\times10^{37}$ times smaller than the GR Lense-Thirring
signal. This is irreversibly unmeasurable.
\label{prop:kerr}
\end{proposition}

\subsection*{2.4 Could a Pulsed High-Field Electromagnet Produce a Detectable
Signal?}

If $|\lam-1|$ were at the drag threshold $1.85\times10^{-2}$:
\begin{equation}
\delta\dot\phi_{\rm DFD}^{|\lam-1|=0.0185}
= 1.85\times10^{-2}\times1.77\times10^{-44}
= 3.3\times10^{-46}~\text{rad/s}.
\end{equation}

Still $5\times10^{31}$ times below Lense-Thirring. \textbf{Conclusion: Kerr
analog frame-dragging is not a viable DFD experimental target at any
accessible coupling strength. The DFD deviation from GR in frame-dragging
is unmeasurable.}

\subsection*{2.5 Qualitative Distinction: When Does DFD Differ from GR?}

DFD and GR predictions for frame-dragging are identical at the level
currently measurable ($\sim1\%$). The qualitative DFD distinction is:
\begin{enumerate}[nosep]
\item A \emph{purely electromagnetic} rotating source (e.g., a rotating
  magnetic toroid with no mass) produces \emph{zero} frame-dragging in GR
  (stress-energy of EM is small) but \emph{non-zero} in DFD (psi-field
  couples to $\mathcal{L}_{\rm EM}$ directly). This is the signature
  to test --- but the magnitude is $\sim 10^{-44}|\lam-1|$~rad/s per unit
  mass equivalent, as computed above.
\item MOND-regime astrophysical systems ($a < a_0$) where the psi-field
  departs from the GR value could show excess frame-dragging. This is
  the P7/P11 connection.
\end{enumerate}

%=============================================================================
\section*{W3i3 \S3: Atom Interferometer Detectability}
\label{sec:aion}
%=============================================================================

\subsection*{3.1 Mach-Zehnder Phase Formula}

For a standard Mach-Zehnder atom interferometer with effective wave
vector $k_{\rm eff} = 2k_L = 4\pi/\lambda_L$ and pulse separation time
$T_{\rm AI}$, the differential phase from an acceleration $a_\psi
= (c^2/2)\nabla\dpsi$:
\begin{equation}
\Delta\phi_{\rm AI} = k_{\rm eff}\,a_\psi\,T_{\rm AI}^2.
\label{eq:AI_phase}
\end{equation}

The minimum detectable acceleration (single shot, shot noise):
\begin{equation}
a_{\min}^{(\rm SS)} = \frac{1}{k_{\rm eff}\,T_{\rm AI}^2\,\sqrt{N_{\rm at}}},
\label{eq:a_min_SS}
\end{equation}
where $N_{\rm at}$ is the atom number per shot.

\subsection*{3.2 AION-100 and MAGIS-100 Specifications}

Both instruments use $^{87}$Sr on the $^1\text{S}_0 \to\,^3\text{P}_0$
clock transition, $\lambda_L = 698$~nm.

\begin{center}
\begin{tabular}{@{}lcc@{}}
\toprule
\textbf{Parameter} & \textbf{Value} & \textbf{Notes} \\
\midrule
$k_{\rm eff}$ & $1.80\times10^7$~m$^{-1}$ & $4\pi/(698\,\text{nm})$ \\
$T_{\rm AI,max}$ & 1.3~s & 100~m baseline \\
$N_{\rm at}$ & $10^6$ atoms/shot & Design spec \\
Shot-noise phase & $10^{-3}$~rad & $1/\sqrt{N_{\rm at}}$ \\
Repetition rate & $\sim1/3$~Hz & $2T_{\rm AI}+t_{\rm dead}$ \\
Daily shots $N_d$ & $\sim 10^4$ & \\
\bottomrule
\end{tabular}
\end{center}

Single-shot acceleration sensitivity:
\begin{equation}
a_{\min}^{(\rm SS)} = \frac{10^{-3}}{1.80\times10^7\times(1.3)^2}
= \frac{10^{-3}}{3.04\times10^7}
= 3.3\times10^{-11}~\text{m/s}^2.
\label{eq:a_min_SS_numerical}
\end{equation}

\textbf{1-day integrated sensitivity}:
\begin{equation}
a_{\min}^{(1d)} = \frac{a_{\min}^{(\rm SS)}}{\sqrt{N_d}}
= \frac{3.3\times10^{-11}}{\sqrt{10^4}}
= 3.3\times10^{-13}~\text{m/s}^2.
\label{eq:a_min_1day}
\end{equation}

\subsection*{3.3 Minimum Detectable Psi-Gradient}

From $a_\psi = (c^2/2)|\nabla\dpsi|$:
\begin{equation}
|\nabla\dpsi|_{\min}^{(1d)}
= \frac{2\,a_{\min}^{(1d)}}{c^2}
= \frac{2\times3.3\times10^{-13}}{(3\times10^8)^2}
= \frac{6.6\times10^{-13}}{9\times10^{16}}
= 7.3\times10^{-30}~\text{m}^{-1}.
\label{eq:dpsi_min}
\end{equation}

\begin{theorem}[Minimum detectable psi-gradient: AION/MAGIS-100, 1 day]
\label{thm:aion}
With 1-day integration at AION-100 or MAGIS-100 design sensitivity:
\begin{equation}
\boxed{
|\nabla\dpsi|_{\min}^{(1d)} = 7.3\times10^{-30}~\text{m}^{-1}.
}
\end{equation}
\end{theorem}

\subsection*{3.4 Required $|\lam-1|$ to Produce This Gradient}

The psi-gradient from a 100~kW EM source at distance $d = 1$~m
(from W3i2 analysis):
\begin{equation}
|\nabla\dpsi|_{\rm EM} \approx |\lam-1|\times1.65\times10^{-39}~\text{m}^{-1}.
\label{eq:gradient_EM_source}
\end{equation}

Setting $|\nabla\dpsi|_{\rm EM} = |\nabla\dpsi|_{\min}^{(1d)}$:
\begin{equation}
|\lam-1|_{\rm AION,req} = \frac{7.3\times10^{-30}}{1.65\times10^{-39}}
= 4.4\times10^9.
\label{eq:lambda_AION}
\end{equation}

\textbf{Gap}: $4.4\times10^9 / 4.9\times10^{-7} = 9\times10^{15}$.

\textbf{Negative finding}: Even with AION/MAGIS-100 sensitivity after
1 day of integration, detecting the psi-gradient from a 100~kW source
at 1~m requires $|\lam-1| = 4.4\times10^9$, nine orders of magnitude
above the drag threshold and $10^{16}$ above the current passive bound.

\subsection*{3.5 Galactic-Center Psi-Gradient vs AION Sensitivity}

For comparison with a natural source: the galactic-center psi-gradient
at Earth's location ($d_{\rm GC} = 8$~kpc $= 2.47\times10^{20}$~m,
$M_{\rm GC} = 4\times10^6\,M_\odot = 8\times10^{36}$~kg):
\begin{equation}
|\nabla\dpsi|_{\rm GC} = \frac{2GM_{\rm GC}}{c^2 d_{\rm GC}^2}
= \frac{2\times6.674\times10^{-11}\times8\times10^{36}}
  {9\times10^{16}\times6.10\times10^{40}}
= \frac{1.069\times10^{27}}{5.49\times10^{57}}
= 1.95\times10^{-31}~\text{m}^{-1}.
\label{eq:GC_gradient}
\end{equation}

Ratio: $|\nabla\dpsi|_{\rm GC}/|\nabla\dpsi|_{\min}^{(1d)}
= 1.95\times10^{-31}/7.3\times10^{-30} = 0.027$.

\textbf{The galactic-center psi-gradient is only 2.7\% of the AION-100
minimum detectable gradient in 1 day.} It would require $\approx 1400$~days
of integration to detect even this natural signal. And this is again the
standard Newtonian gravity signal, not a DFD-specific deviation.

\begin{corollary}[AION/MAGIS cannot detect DFD-specific psi-gradients]
No atom interferometer of AION/MAGIS scale can detect either:
(a) the EM-sourced psi-gradient at any realistic coupling, or
(b) the DFD deviation from GR in naturally occurring psi-gradients.
The only AION/MAGIS signal within reach is standard Newtonian gravity
measurement (already demonstrated) and the search for ultralight dark
matter through psi-field oscillations (not P1 territory).
\label{cor:AION_limitation}
\end{corollary}

\subsection*{3.6 Next-Generation: Space-Based Atom Interferometry}

Proposed space atom interferometers (AEDGE, MAGIS-Space) with $T_{\rm AI}
\sim 300$~s and $L \sim 10^5$~m baseline would achieve:
\begin{equation}
a_{\min}^{\rm space,1yr} \approx \frac{a_{\min}^{(\rm SS,space)}}
{\sqrt{N_{yr}}}
\approx \frac{10^{-3}}{k_{\rm eff}(300)^2\sqrt{3\times10^6}}
= \frac{10^{-3}}{1.80\times10^7\times9\times10^4\times1.73\times10^3}
= \frac{10^{-3}}{2.81\times10^{15}}
= 3.6\times10^{-19}~\text{m/s}^2.
\end{equation}

\textcolor{conjecturecolor}{[CONJECTURE: AEDGE design parameters are
preliminary; $T_{\rm AI} = 300$~s assumes ACES-level microgravity
environment.]}

Corresponding minimum psi-gradient: $8\times10^{-36}$~m$^{-1}$.
Required $|\lam-1|$ for 100~kW source: $\sim5\times10^{3}$. Still
$10^{10}$ above the passive bound.

%=============================================================================
\section*{W3i3 \S4: Non-Reciprocal Drag Effects}
\label{sec:nonreciprocal}
%=============================================================================

\subsection*{4.1 Symmetry Analysis: Is the DFD Drag Force Velocity-Symmetric?}

The standard DFD drag formula (P1 deep analysis, confirmed):
\begin{equation}
F_D = \half\rho_0\,e^{3\dpsi}\,C_D(\Rey_\psi)\,A\,v^2.
\label{eq:drag_standard}
\end{equation}

This is manifestly symmetric in $v \to -v$ \emph{if $\dpsi$ is spatially
uniform}. In a spatially varying $\dpsi$ field, the body traverses
different values of $\dpsi$ depending on which face leads, breaking the
symmetry.

\subsection*{4.2 First Non-Reciprocal Mechanism: Body Straddling a
Gradient}

Consider a body of length $\ell_b$ moving along $x$ through a linear
psi-gradient $G_\psi = \partial\dpsi/\partial x$. The effective
$\dpsi$ experienced by the \emph{leading face} differs from the
\emph{trailing face}:

For $+v$ motion (nose into gradient):
\begin{equation}
\dpsi_{\rm nose}^{(+)} = \dpsi_0 + G_\psi\frac{\ell_b}{2},
\qquad
\dpsi_{\rm tail}^{(+)} = \dpsi_0 - G_\psi\frac{\ell_b}{2}.
\end{equation}

The drag force acts primarily on the nose (form drag from the pressure
distribution):
\begin{equation}
F_{D,+} \approx \half\rho_0\,e^{3(\dpsi_0 + G_\psi\ell_b/2)}\,C_D A\,v^2
\approx F_{D,0}\!\left(1 + \tfrac{3}{2}G_\psi\ell_b\right).
\label{eq:drag_plus}
\end{equation}

For $-v$ motion (tail into gradient; the \emph{same} physical face now leads):
\begin{equation}
F_{D,-} \approx F_{D,0}\!\left(1 + \tfrac{3}{2}G_\psi(-\ell_b)\right)
= F_{D,0}\!\left(1 - \tfrac{3}{2}G_\psi\ell_b\right).
\label{eq:drag_minus}
\end{equation}

\textbf{Non-reciprocal drag asymmetry}:
\begin{equation}
\frac{F_{D,+} - F_{D,-}}{F_{D,0}}
= 3\,G_\psi\,\ell_b
= 3\,\ell_b\,\frac{\partial\dpsi}{\partial x}.
\label{eq:nonrecip_asym}
\end{equation}

\begin{theorem}[Non-Reciprocal Drag Asymmetry, Linear Gradient]
\label{thm:nonrecip}
A body of length $\ell_b$ traversing a linear psi-gradient
$G_\psi = \partial\dpsi/\partial x$ along its axis experiences a
fractional drag asymmetry between $+v$ and $-v$ motion:
\begin{equation}
\boxed{
\frac{\Delta F_D}{F_{D,0}} = 3\,\ell_b\,G_\psi.
}
\end{equation}
The effect is linear in $G_\psi$ and $\ell_b$, and is zero for a
spatially uniform psi-field. This is a \emph{new} P1 prediction
not present in prior iterations.
\end{theorem}

\subsection*{4.3 Second Non-Reciprocal Mechanism: Psi-Body Force}

The DFD body force on a fluid element:
$\mathbf{f}_\psi = \rho_{\rm eff}(c^2/2)\nabla\dpsi$.
For the fluid flowing past a vehicle in the $+x$ direction, this body
force acts as an effective current source that modifies the pressure
distribution asymmetrically. The net force:
\begin{equation}
F_{\rm body,\psi} = \rho_{\rm eff}\,\frac{c^2}{2}\,G_\psi\,V_{\rm body}.
\label{eq:body_force_NR}
\end{equation}

Note: in Earth's gravitational psi-gradient, this reproduces buoyancy
(as expected). The non-trivial non-reciprocal component is the
\emph{change} in this force between $+v$ and $-v$ transit:
\begin{align}
\Delta F_{\rm body,NR}
&\approx \rho_0 e^{3\dpsi}\frac{c^2}{2}G_\psi V_{\rm body}
  \times(e^{3G_\psi\ell_b} - e^{-3G_\psi\ell_b}) \notag\\
&\approx \rho_0 e^{3\dpsi}\frac{c^2}{2}G_\psi V_{\rm body}
  \times 6G_\psi\ell_b.
\label{eq:body_force_asym}
\end{align}

For a 10~m submarine ($V_{\rm body} \sim 1000$~m$^3$) in a man-made
psi-gradient from a 100~kW source at 10~m, $G_\psi \sim |\lam-1|
\times1.65\times10^{-40}$~m$^{-2}$:
\begin{equation}
\Delta F_{\rm body,NR}
\approx 10^3\times 1\times 4.5\times10^{16}
  \times (|\lam-1|\times1.65\times10^{-40})
  \times 1000\times6\times10 \times (|\lam-1|\times1.65\times10^{-40})\times10
\end{equation}
This scales as $|\lam-1|^2$ and is completely negligible even at the
drag threshold.

\subsection*{4.4 Quantitative Assessment of Non-Reciprocal Drag}

\begin{center}
\begin{tabular}{@{}lcc@{}}
\toprule
\textbf{Scenario} & $G_\psi$ (m$^{-1}$) & $\Delta F_D / F_{D,0}$ \\
\midrule
Earth gravity at surface & $3.1\times10^{-16}$ & $9.3\times10^{-15}$ \\
100~kW source, 1~m, bound $|\lam-1|$ & $|\lam-1|\times1.65\times10^{-39}$ & $|\lam-1|\times5\times10^{-38}$ \\
100~kW source, 1~m, drag threshold & $3\times10^{-41}$ & $9\times10^{-40}$ \\
\bottomrule
\end{tabular}
\end{center}

\textbf{Conclusion}: Non-reciprocal drag is a real DFD prediction but at
levels $10^{-39}$--$10^{-14}$ relative to the baseline drag force for any
physically realizable scenario. It is not measurable.

%=============================================================================
\section*{W3i3 \S5: Cross-Patent Synergies}
\label{sec:cross_patent}
%=============================================================================

\subsection*{5.1 P1 + P6: Combined SNR Analysis}

\subsubsection*{5.1.1 P6 Sensor Capabilities at the P1 Operating Point}

The W3i2 P6 update established a 10-node, $K=3$ combined sensitivity of
$3.7\times10^{-24}$~Hz$^{-1/2}$ in strain for decihertz GW. Converting
to acceleration sensitivity at the P1 length scale ($L = 5$~m) and
frequency ($f \sim 1$~Hz for a crossing event):
\begin{equation}
a_{\min}^{(P6)} = \frac{h_{\min}^{(P6)}\,f^2\,L}{2}
= \frac{3.7\times10^{-24}\times1\times5}{2}
= 9.3\times10^{-24}~\text{m/s}^2~\text{Hz}^{-1/2}.
\label{eq:P6_accel_sensitivity}
\end{equation}

This is exceptional sensitivity. Now compute the psi-gradient acceleration
at the P1 drag threshold:
\begin{equation}
a_{\psi}^{(P1\,threshold)}
= \frac{c^2}{2}|\nabla\dpsi|
\approx \frac{c^2}{2}\times\frac{|\dpsi|_{\rm 1\%}}{R_b}
= \frac{9\times10^{16}}{2}
  \times\frac{3.4\times10^{-3}}{5}
= 3.06\times10^{13}~\text{m/s}^2.
\label{eq:P1_threshold_accel}
\end{equation}

The P6 sensor is saturated by $10^{37}$ at the P1 drag threshold. The
P6 sensor is a \emph{null instrument} for P1 drag-reduction physics.

\subsubsection*{5.1.2 The Relevant P6 Application: Passive Psi-Background
Sensing}

The combination P1+P6 is useful only in a context where:
(a)~P6 detects a natural psi-background (astrophysical or geological), and
(b)~P1 provides a drag-based independent confirmation of the same signal.

For a geological psi-gradient $|\nabla\dpsi|_{\rm geo} \sim 10^{-16}$~m$^{-1}$
(from a $10^{10}$~kg underground anomaly at 1~km):
\begin{align}
a_{\rm geo} &= \frac{c^2}{2}\times10^{-16} = 4.5\times10^0 = 4.5~\text{m/s}^2. \\
\Delta F_D/F_D^0 &= e^{3\dpsi_{\rm geo}} - 1 \approx 3|\dpsi|_{\rm geo}.
\end{align}

Even this tiny gradient overwhelms P6 sensitivity. The two instruments
operate in completely non-overlapping regimes.

\begin{proposition}[P1 and P6 are complementary, not additive]
\label{prop:P1P6}
The P1 drag-sensing modality and P6 quantum sensors probe non-overlapping
psi-signal regimes:
\begin{itemize}[nosep]
\item P6 operates at $|\nabla\dpsi| \lesssim 10^{-26}$~m$^{-1}$ (sub-femto-$g$)
\item P1 drag-sensing operates at $|\nabla\dpsi| \gtrsim 10^{-4}$~m$^{-1}$
  (near drag threshold)
\end{itemize}
A combined P1+P6 measurement does not improve SNR over either alone. The
two patents address different rungs of the psi-field amplitude ladder.
\end{proposition}

\subsubsection*{5.1.3 The One Overlapping Application: Kerr-Scale
Astrophysics}

\textcolor{conjecturecolor}{[CONJECTURE]: In the strong-field regime near
a neutron star ($|\dpsi_{\rm NS}| \sim 0.4$), the P6 5470$\times$ combined
sensitivity might detect the DFD post-Newtonian correction to the psi-field
profile. The P1 drag physics in neutron-star atmosphere would then be at
$|\dpsi| \sim 0.4$ --- well above the drag threshold. In this hypothetical
scenario, P1 and P6 would simultaneously activate. However, this is not an
engineerable experimental configuration.}

\subsection*{5.2 P1 + P5: Timing Anomalies from Psi-Gradients}

\subsubsection*{5.2.1 The P5 Nonreciprocal Timing Signal}

The W3i2 P5 update (Modane prediction) gives the DFD nonreciprocal
timing delay for a light path of length $L$ at angle $\theta$ to
$\nabla\dpsi$:
\begin{equation}
\delta t_{\rm NR} = \frac{2L|\nabla\dpsi|\sin\theta}{c}.
\label{eq:P5_NR}
\end{equation}

For the P1 psi-bubble at drag-reduction operating point ($|\dpsi| = 0.3$,
radius $R_b = 5$~m, sharp boundary):
\begin{equation}
|\nabla\dpsi|_{\rm bubble} \approx \frac{|\dpsi|}{R_b} = \frac{0.3}{5} = 0.06~\text{m}^{-1}.
\label{eq:bubble_gradient}
\end{equation}

For a light path $L = 10$~m transversing the bubble at $\theta = 90°$:
\begin{equation}
\delta t_{\rm NR}^{(\rm bubble)} = \frac{2\times10\times0.06}{3\times10^8}
= \frac{1.2}{3\times10^8}
= 4.0~\text{ns}.
\label{eq:bubble_NR}
\end{equation}

\begin{theorem}[P1 psi-bubble produces a measurable P5 timing signal]
\label{thm:P1P5}
At the P1 drag-reduction operating point ($|\dpsi| = 0.3$, $R_b = 5$~m),
the DFD nonreciprocal timing delay through a 10~m transverse optical path
traversing the bubble boundary is:
\begin{equation}
\boxed{
\delta t_{\rm NR}^{(\rm bubble)} = 4.0~\text{ns}.
}
\end{equation}
With atomic-clock precision $\sigma_t = 1$~ps (optical clock, current SOA),
the SNR:
\begin{equation}
\SNR = \frac{4.0~\text{ns}}{1~\text{ps}} = 4000.
\end{equation}
This is the strongest laboratory-scale DFD signal identified across all
15 patents in the research programme.
\end{theorem}

\subsubsection*{5.2.2 Comparison with Modane Baseline}

The W3i2 P5 analysis predicts a Modane nonreciprocal timing signal of
6.2~ps from the galactic psi-gradient over a 4-node, 300~m baseline. The
P1 bubble signal is:
\begin{equation}
\frac{\delta t_{\rm NR}^{(\rm bubble)}}{\delta t_{\rm NR}^{(\rm Modane)}}
= \frac{4.0~\text{ns}}{6.2~\text{ps}} = 645.
\end{equation}

The P1 bubble is 645 times stronger than the Modane cosmic signal --- in a
laboratory at the operator's complete control. The P1 psi-bubble is the
ideal \emph{calibration source} for the P5 timing network.

\subsubsection*{5.2.3 Dense-Object Timing Anomaly: Neutron Star Environment}

For a neutron star with $M_{\rm NS} = 1.4\,M_\odot$, $R_{\rm NS} = 10$~km:
\begin{equation}
|\nabla\dpsi|_{\rm NS,surface}
= \frac{2GM_{\rm NS}}{c^2 R_{\rm NS}^2}
= \frac{2\times6.674\times10^{-11}\times2.8\times10^{30}}
  {9\times10^{16}\times10^8}
= \frac{3.738\times10^{20}}{9\times10^{24}}
= 4.15\times10^{-5}~\text{m}^{-1}.
\label{eq:NS_gradient}
\end{equation}

For a hypothetical 1~km baseline at the NS surface:
\begin{equation}
\delta t_{\rm NR}^{\rm NS} = \frac{2\times10^3\times4.15\times10^{-5}}{3\times10^8}
= 2.8\times10^{-10}~\text{s} = 0.28~\text{ns}.
\label{eq:NS_NR}
\end{equation}

The NS psi-field is in the strong-gravity regime ($|\dpsi_{\rm NS}|
= 2GM/(c^2 R) \approx 0.41$). At this amplitude, the DFD
post-Newtonian deviation from GR is $\order{|\dpsi|^2} \sim 17\%$.
This \emph{could} produce a timing anomaly relative to GR of:
\begin{equation}
\delta t_{\rm DFD-GR}^{\rm NS} \approx 0.17\times\delta t_{\rm NR}^{\rm NS}
= 0.17\times0.28~\text{ns} = 48~\text{ps}.
\label{eq:NS_DFD_GR}
\end{equation}
\textcolor{conjecturecolor}{[CONJECTURE: the 17\% deviation is an
order-of-magnitude estimate; the full strong-field DFD metric
deviates from GR at order $|\dpsi|^2$, but the exact coefficient
requires the DFD post-Newtonian expansion, not yet in the corpus.]}

The 48~ps timing anomaly near a neutron star is detectable with current
optical clock / pulsar timing methods. \textbf{This is a genuine DFD
prediction distinguishable from GR for neutron star environments.}

\subsubsection*{5.2.4 P1+P5 Combined Experimental Design}

The optimal combined experiment uses:
\begin{enumerate}[nosep]
\item A P1 actuator (SRF cavity at $|\lam-1| = 1.85\times10^{-2}$) to
  generate a psi-bubble of radius $R_b = 5$~m.
\item A P5 fiber-linked clock network with 10~m baseline traversing the
  bubble boundary.
\item The 4.0~ns nonreciprocal timing signature confirms the psi-bubble.
\item Progressive reduction of $|\lam-1|$ maps the signal as $\delta t_{\rm NR}
  \propto |\dpsi| \propto |\lam-1|$ (linear in coupling).
\end{enumerate}

This experiment would provide the first \emph{direct laboratory confirmation}
of the DFD psi-field mechanism if $|\lam-1|$ can reach $1.85\times10^{-2}$.

%=============================================================================
\section*{W3i3 \S6: Testable Predictions}
\label{sec:predictions}
%=============================================================================

\noindent The following predictions are new in W3i3:

\begin{enumerate}[label=\textbf{P\arabic*.}, leftmargin=*, nosep, itemsep=4pt]

\item \textbf{Gravitational-Doppler signal is indistinguishable from GR
  for background fields}. The DFD deviation from GR in the clock rate
  for a LEO satellite is $\sim 2\times10^{-19}$ --- five orders of magnitude
  below GPS clock noise. GPS data cannot constrain $|\lam-1|$ beyond what
  the DESY bound already provides. \emph{Negative prediction confirmed
  analytically.}

\item \textbf{Controlled-source psi signal via LEO clock requires optical
  clock array on ISS plus $|\lam-1| > 7\times10^5$.} If future optical
  clock networks (ISS-based) achieve $\sigma_y \sim 10^{-18}$, the
  detectable $|\lam-1|$ threshold drops from $7\times10^8$ (GPS) to
  $7\times10^5$. This remains $10^{12}$ above the passive bound but is
  $10^7$ below the drag threshold. A dedicated search is physically
  motivated but not near-term.

\item \textbf{Non-reciprocal drag asymmetry formula}: $\Delta F_D/F_{D,0}
  = 3\ell_b\,|\nabla\dpsi|$. For the P1 psi-bubble at drag threshold
  ($|\nabla\dpsi| = 0.06$~m$^{-1}$, $\ell_b = 10$~m): asymmetry $= 1.8$.
  An object \emph{inside} a psi-gradient that it itself generates would
  experience an 80\% drag difference between fore and aft motion. This is
  a new internal-consistency check for the P1 psi-bubble design.

\item \textbf{P1 psi-bubble: 4.0~ns nonreciprocal timing signal at
  drag threshold.} Detectable with SNR 4000 using optical clocks. This is
  the calibration signal for P5 deployment. Testable with existing
  technology \emph{if} $|\lam-1|$ reaches $1.85\times10^{-2}$.

\item \textbf{DFD post-Newtonian deviation near neutron stars: 48~ps
  timing anomaly.} At the neutron star surface ($|\dpsi| \approx 0.41$),
  the DFD post-Newtonian correction to the Shapiro/Lense-Thirring timing
  is $\sim$17\%, producing a 48~ps anomaly detectable by pulsar timing
  arrays \textcolor{conjecturecolor}{[CONJECTURE: coefficient uncertain;
  requires DFD strong-field PPN calculation]}.

\item \textbf{Kerr-analog frame-dragging: DFD correction is $10^{38}$
  below GR prediction.} Completely unmeasurable at any coupling strength
  below the physical bound $|\lam-1| \lesssim 1$. This closes the Kerr
  analog question conclusively.

\item \textbf{AION/MAGIS-100 cannot detect DFD-specific psi-gradients}
  from any controlled source. Required coupling is $|\lam-1| = 4.4\times10^9$
  for a 100~kW source at 1~m in 1 day. The gap to the passive bound is
  $10^{16}$. Atom interferometers probe standard gravity, not DFD-specific
  deviations.

\end{enumerate}

%=============================================================================
\section*{W3i3 Summary Tables}
%=============================================================================

\begin{center}
\textbf{Table 1: New Derivations in W3i3 (not in W3i1 or W3i2)}
\vspace{0.5em}

\begin{tabular}{@{}p{5cm}cc@{}}
\toprule
\textbf{Result} & \textbf{Value} & \textbf{Verdict} \\
\midrule
LEO clock DFD vs GR difference & $2\times10^{-19}$ & Negative \\
LEO Eccentric orbit DFD excursion & $1.32\times10^{-10}$ & GR-equivalent \\
Controlled source, GPS detectability & $|\lam-1| > 7\times10^8$ & Negative \\
DFD Kerr-analog precession & $8.7\times10^{-52}$~rad/s & Negative \\
AION-100, 1-day $|\nabla\dpsi|_{\min}$ & $7.3\times10^{-30}$~m$^{-1}$ & Derived \\
Required $|\lam-1|$ for AION & $4.4\times10^9$ & Negative \\
Non-reciprocal drag formula & $3\ell_b|\nabla\dpsi|$ & New theorem \\
NR drag at drag threshold & 1.8 (180\% aft vs fore) & New prediction \\
P1+P6 SNR gain & None (P6 saturated) & Negative \\
P1 bubble timing signal & 4.0~ns, SNR=4000 & \textbf{Key result} \\
P1 vs Modane signal ratio & 645$\times$ & \textbf{Key result} \\
NS timing anomaly (DFD vs GR) & $\sim$48~ps & Conjecture \\
\bottomrule
\end{tabular}
\end{center}

\vspace{0.8em}

\begin{center}
\textbf{Table 2: Complete Technology Threshold Hierarchy (W3i1 through W3i3)}
\vspace{0.5em}

\begin{tabular}{@{}p{5.5cm}ccc@{}}
\toprule
\textbf{Effect} & $|\dpsi|$ & $|\lam-1|$ & Gap \\
\midrule
K-H suppression (sub, 12~m/s) & $4.9\times10^{-9}$ & $\sim5\times10^{-4}$ & $10^3$ \\
K-H suppression (aircraft, 300~m/s) & $1.2\times10^{-7}$ & $\sim10^{-2}$ & $2\times10^4$ \\
1\% drag reduction (bluff body) & $3.4\times10^{-3}$ & $1.85\times10^{-2}$ & $3.8\times10^4$ \\
P5 timing calibration (P1 bubble) & $0.3$ & $1.85\times10^{-2}$ & $3.8\times10^4$ \\
10~dB acoustic cloaking & $0.253$ & $\sim1$ & $\gg1$ \\
60\% drag reduction & $0.305$ & $\sim1$ & $\gg1$ \\
10~dB EM cloaking & $0.632$ & $\gg1$ & $\gg1$ \\
\midrule
\multicolumn{4}{l}{\textit{Passive bound}: $|\lam-1| < 4.9\times10^{-7}$ (DESY, $\mu_0$-corrected).} \\
\bottomrule
\end{tabular}
\end{center}

%=============================================================================
\section*{Open Questions for Iteration 4}
%=============================================================================

\begin{enumerate}[nosep, itemsep=4pt]

\item \textbf{Full DFD post-Newtonian expansion}: Derive the PPN parameters
  $\beta$, $\gamma$ for DFD and compute the post-Newtonian clock rate to
  order $|\dpsi|^2$. Confirm or correct the $\sim$17\% NS timing anomaly
  estimate.

\item \textbf{Non-reciprocal drag inside the psi-bubble}: The 180\%
  forward/aft asymmetry predicted by Prediction P3 above assumes the body
  is immersed \emph{inside} the gradient. Model the actual bubble geometry:
  what is the asymmetric drag on an object crossing from outside to inside
  the bubble, vs purely tangential motion?

\item \textbf{P5 timing calibration protocol}: Design a concrete experiment:
  source (SRF cavity at 35~GHz), propagation path (10~m fiber), clock
  (optical Sr lattice). What are the dominant systematics (e.g., thermal
  expansion, Sagnac, magnetic field gradients) and how do they compare
  to the 4.0~ns target signal?

\item \textbf{Neutron star pulsar timing as DFD test}: Pulsar timing arrays
  (NANOGrav, PPTA) have $\sim$100~ns timing precision over multi-year
  baselines. A 48~ps anomaly in single-pulse timing would require $\sim
  10^{-4}$ relative accuracy. Identify the most favourable pulsar (closest,
  most stable, deepest potential well) and compute the expected DFD timing
  deviation.

\item \textbf{Strong-field DFD cloaking}: At $|\dpsi| \sim 0.4$ (neutron
  star surface analog, achievable in principle at drag threshold for a
  dense medium), the acoustic impedance mismatch is $e^{5\times0.4/2}
  = e = 2.71$, giving $r_{\rm ac} = (e-1)/(e+1) = 0.46$. This is 46\%
  acoustic reflection --- far above any cloaking operating point. Does P1
  cloaking have a maximum at $|\dpsi| \sim 0.25$ and degrade at higher
  coupling? Derive the acoustic signature curve as a function of $|\dpsi|$
  from 0 to 1.

\end{enumerate}

\end{document}
